

\section{Simple types with fixed points: system $\muLJ$}
\label{sec:muLJ-prelims}
In this section we recall the system $\muLJ$ from \cite{Clairambault09,Clairambault13interleaving}.
More precisely, we present the `strong' version of $\muLJ$ from \cite{Clairambault13interleaving}.
% Since much development of the system already exists in those earlier papers we shall be rather brief.



\subsection{The sequent calculus $\muLJ$}

% \begin{defn}[Pretypes and free variables]
\emph{Pretypes}, written $\sigma,\tau$ etc., are generated by the following grammar:
\[
\sigma,\tau \ \bnf\  X \ \vert \ \unit \ \vert \ \sigma + \tau \ \vert \ \sigma \times \tau \ \vert \ \sigma \arrow  \tau  \ \vert \ \mu X \sigma \ \vert \ \nu X \sigma      
\]


\emph{Free (second-order) variables} of a pretype are defined as expected, construing $\mu$ and $\nu$ as binders:
\begin{itemize}
    \item $\fv X \dfn \{X\}$
    \item $\fv \unit \dfn \emptyset$
    \item $\fv {\sigma \star \tau} \dfn \fv \sigma \cup \fv \tau$, for $\star \in \{+,\times,\arrow\}$
    \item $\fv{\kappa X \, \sigma} \dfn \fv \sigma \setminus \{X\}$, for $\kappa \in \{\mu,\nu\}$
\end{itemize}
A pretype is \emph{closed} if it has no free variables (otherwise it is \emph{open}).
% \end{defn}

Throughout this work we shall assume some standard conventions on variable binding, in particular that each occurrence of a binder $\mu$ and $\nu$ binds a variable distinct from all other binder occurrences in consideration.
This avoids having to deal with variable renaming explicitly.
We follow usual bracketing conventions, in particular writing, say, $\rho \arrow \sigma\arrow \tau$ for $(\rho \arrow (\sigma \arrow \tau))$.
Binders $\mu X $ and $\nu Y$ bind as strongly as possible but we may write, say, $\mu X, Y .\,  \sigma \to \tau$ for $\mu X \mu Y (\sigma \to \tau)$. 


\begin{defn}
    [Types and polarity]
    \emph{Positive} and \emph{negative} variables in a pretype are defined as expected:
\begin{itemize}     
\item $X$ is positive in $X$.
        \item $1$ is positive and negative in $X$.
        \item if $\sigma,\tau$ are positive (negative) in $X$ then so is $\sigma \star \tau$, for $\star \in \{+,\times\}$.
        \item if $\sigma$ is negative (positive) in $X$ and $\tau $ is positive (resp., negative) in $X$, then $\sigma \to \tau$ is positive (resp., negative) in $X$.
        \item if $\sigma$ is positive (negative) in $X$ then so is $\kappa Y \sigma$ (resp.), for $\kappa \in \{\mu,\nu\}$, both when $Y=X$ and $Y \neq X$.
    \end{itemize}
A pretype is a \emph{type} (or even \emph{formula}) if, for any subexpression $\kappa X \sigma$, $\sigma$ is positive in $X$.
The notions of \emph{(type) context} and \emph{substitution} are defined as usual. 
% In particular we write $\sigma[\tau/X]$ for the type obtained by replacing every instance of $X$ in $\sigma$ by $\tau$.\anupam{do we need this?}
\end{defn}

\begin{remark}
    [Positivity vs strict positivity]
    Many authors require variables bound by fixed point operators to appear in \emph{strictly} positive position, i.e.\ never under the left-scope of $\arrow$.
    Like Clairambault \cite{Clairambault09,Clairambault13interleaving} we do not impose this stronger requirement, requiring only positivity in the usual syntactic sense.
\end{remark}

\begin{defn}[System $\muLJ$]
A \emph{cedent}, written $\Sigma, \Gamma$ etc., is just a list of types. 
A \emph{sequent} is an expression $\Sigma \seqar \sigma$.
The symbol $\seqar$ is, formally, just a syntactic delimiter (but the arrow notation is suggestive).
The system $\muLJ$ is given by the rules of Figures~\ref{fig:sequent-calculus-mulj}, \ref{fig:unfolding-rules} and \ref{fig:co-iteration-rules} (colours may be ignored for now).
The notions of \emph{derivation} (or \emph{proof}) are defined as usual. 
We write $P: \Gamma \seqar \tau$ if $P$ is a derivation of the sequent $\Gamma \seqar \tau$. 
\end{defn}

\begin{remark}
    [General identity and substitutions]
    \label{rem:gen-identity+substitution}
    Note that $\muLJ$ is equipped with a general identity rule, not only for atomic types. 
    This has the apparently simple but useful consequence that typing derivations are closed under substitution of types for free variables, i.e.\ if $\der(X) : \Gamma(X) \seqar \sigma(X)$ in $\muLJ $ (with all occurrences of $X$ indicated), then also $\der(\tau) : \Gamma(\tau) \seqar \sigma(\tau)$ in $\muLJ$ for any type $\tau$.
    Later, this will allow us to derive inductively general functors for fixed points in $\muLJ$ rather than including them natively; this will in turn become important later for verifying our realisability model for $\muLJ$ \anupam{come back to this later, need to mention eta}.
\end{remark}

\begin{remark}
    [$\muLJ$ as a fragment of second-order logic]
    \label{rem:mulj-frag-F}
    We may regard $\muLJ$ properly as a fragment of Girard-Reynolds System $\F$~\cite{girard1972systemF,Reynolds74systemF}, an extension of simple types to a second-order setting. 
    In particular, (co)inductive types may be identified with second-order formulas by:
    \[
    \begin{array}{r@{\ = \ }l}
         \mu X \sigma & \forall X ((\sigma \arrow X) \arrow X) \\
         \nu X \sigma & \exists X (X \times (X \arrow \sigma))
    \end{array}
    \]
    The rules for fixed points in $\muLJ$ are essentially inherited from this encoding, modulo some constraints on proof search strategy.
    % Indeed \emph{cut-elimination} has been obtained directly for fixed point systems by Baelde and Miller~\cite{BaeldeMiller07}.\footnote{Note that their result is stated over (multiplicative additive) linear logic.}
% 
    Later we shall use a different encoding of fixed point types into a second-order setting, namely in arithmetic, as bona fide fixed points, in order to better control logical complexity.\anupam{could say more}
\end{remark}

In proofs that follow, we shall frequently only consider the cases of \emph{least} fixed points ($\mu$-types) and not \emph{greatest} fixed points ($\nu$-types), 
% due to space constraints, 
appealing to `duality' for the latter.
The cases for $\nu$ should be deemed analogous.
% , though we provide omitted cases in appendices nonetheless \todo{make sure this is done}.
% We shall somewhat make up for this by including examples for $\nu$-types throughout \anupam{do this}.
As we shall soon see, in Subsection~\ref{subsec:reduction-to-neg-frag}, we can indeed reduce our consideration to $\nu$-free types, without loss of generality in terms of representable functions.

\begin{remark}
    [Why sequent calculus?]
    Using a sequent calculus as our underlying type system is by no means the only choice. 
    % Other approaches to (co)inductive types exist in the literature too, e.g.\ \todo{cites}.
    However, since we shall soon consider non-wellfounded and circular typing derivations, it is important to have access to a well behaved notion of \emph{formula ancestry}, in order to properly define the usual totality criterion that underlies them.
    This is why the sequent calculus is the formalism of choice in circular proof theory.
\end{remark}


\begin{remark}
    [Variations of the fixed point rules]
    \label{rem:iter-variations}
    It is common to consider context-free and `weak' specialisations of the fixed point rules, e.g.:
    \begin{equation}
    \small
        \label{eq:(co)iterators-as-rules}
            \vlinf{\iter{}}{}{\mu X \sigma(X)\seqar \tau}{\sigma(\tau)\seqar \tau}
    \quad
    \vlinf{\coiter}{}{\rho\seqar \nu X \sigma(X)}{\rho\seqar \sigma(\rho)}
    \end{equation}
    % Here we have omitted identity premisses.
    In the presence of cut the `(co)iterator' rules above are equivalent to those of $\muLJ$ (see \Cref{app:simple-types-with-fixed-points} for some further remarks).
    However since the computational model we presume is cut-reduction, as we shall soon see, it is not appropriate to take them as first-class citizens.
    When giving a semantics that interprets $\cut$ directly, e.g.\ as we do for the term calculi in Section~\ref{sec:term-calculi}, it is often simpler to work with the (co)iterators above.
% 
    In the remainder of this work
    we shall freely use the versions above in proofs too.
\end{remark}





\begin{figure}
    \centering
    \[
    \small
%     \begin{array}{c}
%          \vlinf{\ax}{}{\sigma \seqar \sigma}{}
%    \qquad  
%    \vlinf{\ex}{}{\mygreen{\Gamma}, \red{\tau}, \blue{\sigma}, \orange{\Delta} \seqar \cyan{\gamma}}{\mygreen{\Gamma},\blue{\sigma}, \red{\tau}, \orange{\Delta} \seqar \cyan{\gamma}}
%    \qquad
%    \vliinf{\cut}{}{\mygreen{\Gamma}, \orange{\Delta} \seqar \cyan{\tau}}{\mygreen{\Gamma} \seqar \sigma}{\orange{\Delta}, \sigma \seqar \cyan{\tau}}
%    \\ \\
%    \vlinf{\wk}{}{\mygreen{\Gamma}, \sigma \seqar \cyan{\tau}}{\Gamma \seqar \cyan{\tau}}
%    \qquad
%    \vlinf{\contr}{}{\mygreen{\Gamma}, \red{\sigma} \seqar \cyan{\tau}}{\mygreen{\Gamma}, \red{\sigma}, \red{\sigma} \seqar \cyan{\tau}}
% %   \qquad 
% %   \vlinf{\zero}{}{\Gamma, \zero \seqar \sigma }{}
%    \qquad 
%    \vlinf{\rr \unit}{}{\seqar \unit}{}
%    \qquad 
%    \vlinf{\lr \unit}{}{\mygreen{\Gamma}, \unit \seqar \cyan{\sigma}}{\mygreen{\Gamma} \seqar \cyan{\sigma}}
%    \\ \\
%     \vlinf{\rr \arrow}{}{\mygreen{\Gamma} \seqar \red{\sigma} \arrow \cyan{\tau}}{\mygreen{\Gamma}, \red{\sigma} \seqar \cyan{\tau}}
%    \qquad 
%    \vliinf{\lr \arrow}{}{\mygreen{\Gamma},  \orange{\Delta}, \red{\sigma} \arrow \blue{\tau} \seqar \cyan{\gamma}}{\mygreen{\Gamma} \seqar \red{\sigma}}{\orange{\Delta}, \blue{\tau} \seqar \cyan{\gamma}}
%    \\ \\
%    \vliinf{\rr \times}{}{\mygreen{\Gamma} , \orange{\Delta} \seqar \red{\sigma} \times \blue{\tau}}{\mygreen{\Gamma} \seqar \red{\sigma}}{\orange{\Delta} \seqar \blue{\tau}}
%    \qquad 
%    \vlinf{\lr \times}{}{\mygreen{\Gamma}, \red{\sigma} \times \blue{\tau} \seqar \cyan{\gamma}}{\mygreen{\Gamma},\red{\sigma},\blue{\tau} \seqar \cyan{\gamma}}\\ \\
%    \vlinf{\choice 0 {\rr +}}{}{\mygreen{\Gamma} \seqar \blue{\tau_0} + \tau_1}{\mygreen{\Gamma} \seqar \blue{\tau_0}}\qquad
%     \vlinf{\choice 1 {\rr +}}{}{\mygreen{\Gamma} \seqar \tau_0 + \blue{\tau_1}}{\mygreen{\Gamma} \seqar \blue{\tau_1}}
%    % \qquad
%    %  \vlinf{\rr +^1}{}{\Gamma \seqar \tau_0 + \tau_1}{\Gamma \seqar \tau_1}
%     \qquad 
%     \vliinf{\lr +}{}{\mygreen{\Gamma}, \red{\sigma} +\blue{\tau} \seqar \cyan{\gamma}}{\mygreen{\Gamma}, \red{\sigma} \seqar \cyan{\gamma}}{\mygreen{\Gamma}, \blue{\tau} \seqar \cyan{\gamma}}
%     \end{array}
    \begin{array}{c}
         \vlinf{\ax}{}{\sigma \seqar \sigma}{}
   \qquad  
   \vlinf{\ex}{}{\mygreen{\Gamma}, \red{\tau}, \blue{\sigma}, \orange{\Delta} \seqar \cyan{\gamma}}{\mygreen{\Gamma},\blue{\sigma}, \red{\tau}, \orange{\Delta} \seqar \cyan{\gamma}}
   \qquad
   \vliinf{\cut}{}{\mygreen{\Gamma}, \orange{\Delta} \seqar \cyan{\tau}}{\mygreen{\Gamma} \seqar \sigma}{\orange{\Delta}, \sigma \seqar \cyan{\tau}}
   \\ \\
   \vlinf{\wk}{}{\mygreen{\Gamma}, \sigma \seqar \cyan{\tau}}{\Gamma \seqar \cyan{\tau}}
   \qquad
   \vlinf{\contr}{}{\mygreen{\Gamma}, \red{\sigma} \seqar \cyan{\tau}}{\mygreen{\Gamma}, \red{\sigma}, \red{\sigma} \seqar \cyan{\tau}}
%   \qquad 
%   \vlinf{\zero}{}{\Gamma, \zero \seqar \sigma }{}
   \qquad 
   \vlinf{\rr \unit}{}{\seqar \unit}{}
   \qquad 
   \vlinf{\lr \unit}{}{\mygreen{\Gamma}, \unit \seqar \cyan{\sigma}}{\mygreen{\Gamma} \seqar \cyan{\sigma}}
   \\ \\
    \vlinf{\rr \arrow}{}{\mygreen{\Gamma} \seqar \red{\sigma \arrow \tau}}{\mygreen{\Gamma}, \red{\sigma} \seqar \red{\tau}}
   \qquad 
   \vliinf{\lr \arrow}{}{\mygreen{\Gamma},  \orange{\Delta}, \red{\sigma \arrow \tau} \seqar \cyan{\gamma}}{\mygreen{\Gamma} \seqar \red{\sigma}}{\orange{\Delta}, \red{\tau} \seqar \cyan{\gamma}}
   \\ \\
   \vliinf{\rr \times}{}{\mygreen{\Gamma} , \orange{\Delta} \seqar \red{\sigma \times \tau}}{\mygreen{\Gamma} \seqar \red{\sigma}}{\orange{\Delta} \seqar \red{\tau}}
   \qquad 
   \vlinf{\lr \times}{}{\mygreen{\Gamma}, \red{\sigma \times \tau} \seqar \cyan{\gamma}}{\mygreen{\Gamma},\red{\sigma},\red{\tau} \seqar \cyan{\gamma}}\\ \\
   \vlinf{\choice 0 {\rr +}}{}{\mygreen{\Gamma} \seqar \blue{\tau_0 + \tau_1}}{\mygreen{\Gamma} \seqar \blue{\tau_0}}\qquad
    \vlinf{\choice 1 {\rr +}}{}{\mygreen{\Gamma} \seqar \blue{\tau_0 + \tau_1}}{\mygreen{\Gamma} \seqar \blue{\tau_1}}
   % \qquad
   %  \vlinf{\rr +^1}{}{\Gamma \seqar \tau_0 + \tau_1}{\Gamma \seqar \tau_1}
    \qquad 
    \vliinf{\lr +}{}{\mygreen{\Gamma}, \red{\sigma + \tau} \seqar \cyan{\gamma}}{\mygreen{\Gamma}, \red{\sigma} \seqar \cyan{\gamma}}{\mygreen{\Gamma}, \red{\tau} \seqar \cyan{\gamma}}
    \end{array}
    \]
    \caption{Sequent calculus rules for $\LJ$.}
    \label{fig:sequent-calculus-mulj}
\end{figure}

\anupam{27-02-23: I corrected the colouring here, it is still incorrect in the conference version.}

\begin{figure}
 \[
 \small
   \vlinf{\rr \mu}{}{\mygreen{\Gamma}\seqar  \red{\mu X \sigma(X)}}{\mygreen{\Gamma} \seqar \red{\sigma(\mu X \sigma(X))}}
   \qquad
   \vlinf{\lr \nu }{}{\mygreen{\Gamma},  \red{\nu X\sigma(X)} \seqar \cyan{\tau}}{\mygreen{\Gamma} , \red{\sigma(\nu X \sigma(X))}\seqar \cyan{\tau}}
   \]
    \caption{Some unfolding rules for $\mu$ and $\nu$.}
    \label{fig:unfolding-rules}
\end{figure}

\begin{figure}
    \[
    \small
    \vliinf{\lr \mu}{}{\Gamma,\Delta, \mu X \sigma(X)  \seqar \tau}{\Gamma, \sigma(\rho)\seqar \rho}{\Delta, \rho \seqar \tau}
  \qquad
   \vliinf{\rr \nu }{}{\Gamma , \Delta \seqar \nu X\sigma(X)}{\Gamma \seqar  \tau }{  \Delta, \tau \seqar \sigma(\tau)}
    \]
    \caption{`(Co)iteration' rules for $\mu$ and $\nu$.}
    \label{fig:co-iteration-rules}
\end{figure}


\begin{definition}
    [Functors]
    \label{def:functors-in-muLJ}
    % Let $\sigma(X) $ and $\rho(X)$ be (possibly open) types that are positive (and negative, respectively) in $X$. 
    % For a proof $P:\Gamma, \tau \seqar \tau'$ we define $\sigma(P): \Gamma, \sigma(\tau)\seqar \sigma(\tau')$ and $\rho(P): \Gamma, \rho(\tau')\seqar \rho(\tau)$ by simultaneous induction. If $\sigma$ and/or $\rho$ are formed from smaller types by a connective of $\LJ$, then the corresponding functors are defined as usual. 
    % If $\sigma(X) = \mu Y \sigma'(X,Y)$ then $\sigma(P)$ is,
    %     \[
    %     \small
    %     \vlderivation{
    %     \vlin{\lr \mu}{}{\Gamma, \sigma(\tau) \seqar \sigma(\tau')}{
    %     \vlin{\rr \mu}{}{\Gamma, \sigma'(\tau, \sigma(\tau'))\seqar \sigma(\tau')}{
    %     \vliq{\sigma'(P,\sigma(\tau'))}{}{\Gamma, \sigma'(\tau, \sigma(\tau')) \seqar \sigma'(\tau', \sigma(\tau')) }{
    %     \vlhy{\Gamma, \tau \seqar \tau'}
    %     }
    %     }
    %     }
    %     }
    %     \]
    %     where $\sigma'(P,\sigma(\tau'))$ is obtained from the inductive hypothesis for $\sigma'(P,Y)$ under substitution of $\sigma(\tau')$ for $Y$.
    %     Dually for $\rho(P)$, and similarly for greatest fixed points. 
 Let $\sigma(X) $ and $\rho(X)$ be (possibly open) types that are positive (and negative, respectively) in $X$. 
    For a proof $P:\Gamma, \tau \seqar \tau'$ we define $\sigma(P): \Gamma, \sigma(\tau)\seqar \sigma(\tau')$ and $\rho(P): \Gamma, \rho(\tau')\seqar \rho(\tau)$ by simultaneous induction as follows:
    \begin{itemize}
    \item If $\sigma(X)=X$ then $\sigma(\der)$ is just $\der$. Notice that it is never the  case that $\rho=X$, as $X$ can only occur negatively in $\rho$.
    \item If $\sigma$ and  $\rho$ are $\unit$ or some $Y\neq X$ then $\sigma(\der)$ and ${\tau}(\der)$ are defined respectively as follows: 
    \[
    \small
    \vlderivation{
    \vliq{\wk}{}{\Gamma, \sigma \seqar \sigma}{\vlin{\ax}{}{\sigma \seqar \sigma}{\vlhy{}}}
    }
    \qquad 
    \vlderivation{
    \vliq{\wk}{}{\Gamma, \rho \seqar \rho}{\vlin{\ax}{}{\rho \seqar \rho}{\vlhy{}}}
    }
    \]
    \item If $\sigma=\sigma_1 \imp \sigma_2$ and $\rho= \rho \imp \rho$ then we define $\sigma(\der)$ and $\rho(\der)$ respectively as follows:
    \[
    \footnotesize
    \vlderivation{
    \vlin{\rr \imp}{}{\Gamma, \sigma(\tau)\seqar \sigma(\tau')  }
     {
     \vliq{\contr}{}{\Gamma, \sigma(\tau) , \sigma_1(\tau')\seqar   \sigma_2(\tau')}{
     \vliin{\lr \imp}{}{\Gamma ,\Gamma, \sigma(\tau) , \sigma_1(\tau') \seqar   \sigma_2(\tau')}
       {
      \vldr{\sigma_1(\der)}{\Gamma, \sigma_1(\tau') \seqar \sigma_1(\tau)}
       }
       {
        \vldr{\sigma_2(\der)}{\Gamma, \sigma_2(\tau) \seqar \sigma_2(\tau')}
       }
       }
     }
    } 
     \qquad 
   \vlderivation{
    \vlin{\rr \imp}{}{\Gamma, \rho(\tau')\seqar \rho(\tau)  }
     {
     \vliq{\contr}{}{\Gamma, \rho(\tau'), \rho_1(\tau)\seqar   \rho_2(\tau)}{
     \vliin{\lr \imp}{}{\Gamma, \Gamma, \rho(\tau') , \rho_1(\tau) \seqar   \rho_2(\tau)}
       {
       \vldr{\rho_1(\der)}{\Gamma, \rho_1(\tau) \seqar \rho_1(\tau')}
       }
       {
        \vldr{\rho_2(\der)}{\Gamma, \rho_2(\tau') \seqar \rho_2(\tau)}
       }
       }
     }
    }
    \]
    \item If $\sigma=\sigma_1 \times \sigma_2$ and $\rho= \rho_1 \times \rho_2$  then we define $\sigma(\der)$ and $\rho(\der)$ respectively as follows:
     \[
     \footnotesize
   \vlderivation{
     \vliq{\contr}{}{\Gamma, \sigma(\tau)\seqar \sigma(\tau')  }
      {
      \vliin{\rr \times}{}{\Gamma, \Gamma, \sigma(\tau)\seqar\sigma(\tau)  }
       {
       \vlin{\lr \times }{}{\Gamma,\sigma(\tau)\seqar \sigma_1(\tau')}
         {
         \vldr{\sigma_1(\der)}{\Gamma, \sigma_1(\tau) \seqar \sigma_1(\tau')}
         }
       }
       {
         \vlin{\lr \times }{}{\Gamma,\sigma(\tau)\seqar \sigma_2(\tau')}
          {
             \vldr{\sigma_2(\der)}{\Gamma,  \sigma_2(\tau)\seqar \sigma_2(\tau')}
          }
       }
      }
    } 
 \qquad 
    \vlderivation{
     \vliq{\contr}{}{\Gamma, \rho(\tau')\seqar \rho(\tau)  }
      {
      \vliin{\rr \times}{}{\Gamma, \Gamma, \rho(\tau') \seqar \rho(\tau)   }
       {
       \vlin{\lr \times }{}{\Gamma, \rho(\tau') \seqar \rho_1(\tau)}
         {
         \vldr{\rho_1(\der)}{\Gamma, \rho_1(\tau') \seqar \rho_1(\tau)}
         }
       }
       {
         \vlin{\lr \times }{}{\Gamma, \rho(\tau') \seqar \rho_2(\tau)}
          {
             \vldr{\rho_2(\der)}{\Gamma,  \rho_2(\tau')\seqar \rho_2(\tau)}
          }
       }
      }
    } 
    \]
    \item If $\sigma=\sigma_1 + \sigma_2$ and $\rho= \rho_1 + \rho_2$ then we define $\sigma(\der)$ and $\rho(\der)$ respectively as follows:
    \[
    \footnotesize
   \vlderivation{
    \vliin{\lr +}{}{\Gamma, \sigma(\tau)\seqar  \sigma(\tau') }
     {
     \vlin{\rr +}{}{\Gamma, \sigma_1(\tau) \seqar  \sigma(\tau')}
       {
       \vldr{\sigma_1(\der)}{\Gamma, \sigma_1(\tau) \seqar \sigma_1(\tau')}
       }
     }
     {
      \vlin{\rr +}{}{\Gamma, \sigma_2(\tau) \seqar  \sigma(\tau')}
       {
        \vldr{\sigma_2(\der)}{\Gamma, \sigma_2(\tau) \seqar \sigma_2(\tau')}
       }
     }
    }
    \qquad 
     \vlderivation{
    \vliin{\lr +}{}{\Gamma, \rho(\tau') \seqar \rho(\tau)  }
     {
     \vlin{\rr +}{}{\Gamma, \rho_1(\tau') \seqar \rho(\tau) }
       {
       \vldr{\rho_1(\der)}{\Gamma, \rho_1(\tau') \seqar \rho_1(\tau)}
       }
     }
     {
      \vlin{\rr +}{}{\Gamma, \rho_2(\tau') \seqar \rho(\tau) }
       {
        \vldr{\rho_2(\der)}{\Gamma, \rho_2(\tau') \seqar \rho_2(\tau)}
       }
     }
    }
    \]
\item if $\sigma(X) = \mu Y \sigma'(X,Y)$ and $\rho(X) = \mu Y \rho'(X,Y)$  then we define $\sigma(\der)$ and $\rho(\der)$ respectively as follows:
        \[
        \small
        \vlderivation{
        \vlin{\lr \mu}{}{\Gamma, \sigma(\tau) \seqar \sigma(\tau')}{
        \vlin{\rr \mu}{}{\Gamma, \sigma'(\tau, \sigma(\tau'))\seqar \sigma(\tau')}{
        \vldr{\colorbox{white}{$\scriptstyle\sigma'(P,\sigma(\tau'))$}}{\Gamma, \sigma'(\tau, \sigma(\tau')) \seqar \sigma'(\tau', \sigma(\tau'))}
        }
        }
        }
        \qquad         \vlderivation{
        \vlin{\lr \mu}{}{\Gamma, \rho(\tau') \seqar \rho(\tau)}{
        \vlin{\rr \mu}{}{\Gamma, \rho'(\tau', \rho(\tau))\seqar \rho(\tau)}{
        \vldr{\colorbox{white}{$\scriptstyle\rho'(P,\rho(\tau))$}}{\Gamma, \rho'(\tau', \rho(\tau)) \seqar \rho'(\tau, \rho(\tau))}
        }
        }
        }
        \]
        where $\sigma'(P,\sigma(\tau'))$ (resp., $\rho'(P,\rho(\tau))$) are obtained from the IH for $\sigma'(P,Y)$ (resp., $\rho'(P,Y)$) under substitution of $\sigma(\tau')$ for $Y$ (resp., $\rho(\tau)$), cf.~Remark~\ref{rem:gen-identity+substitution}.
        \item   if $\sigma(X) = \nu Y \sigma'(X,Y)$ and $\rho(X) = \nu Y \rho'(X,Y)$  then we define $\sigma(\der)$ and $\rho(\der)$ respectively as follows: 
            \[
        \small
        \vlderivation{
        \vlin{\rr \nu}{}{\Gamma, \sigma(\tau) \seqar \sigma(\tau')}{
        \vlin{\lr \nu}{}{\Gamma, \sigma(\tau) \seqar \sigma'(\tau', \sigma(\tau))}{
        \vldr{\colorbox{white}{$\scriptstyle\sigma'(P,\sigma(\tau))$}}{\Gamma, \sigma'(\tau, \sigma(\tau)) \seqar \sigma'(\tau', \sigma(\tau))}
        }
        }
        }
        \qquad 
        \vlderivation{
        \vlin{\rr \nu}{}{\Gamma, \rho(\tau') \seqar \rho(\tau)}{
        \vlin{\lr \nu}{}{\Gamma, \rho(\tau')\seqar \rho'(\tau, \rho(\tau'))}{
        \vldr{\colorbox{white}{$\scriptstyle\rho'(P,\rho(\tau'))$}}{\Gamma, \rho'(\tau', \rho(\tau')) \seqar \rho'(\tau, \rho(\tau'))}
        }
        }
        }
        \]
        where $\sigma'(P,\sigma(\tau))$ (resp., $\rho'(P,\rho(\tau'))$) are obtained from the IH for $\sigma'(P,Y)$ (resp., $\rho'(P,Y)$) under substitution of $\sigma(\tau)$ for $Y$ (resp., $\rho(\tau')$), cf.~Remark~\ref{rem:gen-identity+substitution}.
        \end{itemize}
        % Further details are in Appendix~\ref{subsec:functoriality-in-muLJ}.
    %     \item if $\sigma$ and/or $\rho$ are formed from smaller types by a connective of $\LJ$, then the corresponding functors are defined as usual.
    %     \item if $\sigma(X) = \mu Y \sigma'(X,Y)$ then $\sigma(P)$ is,
    %     \[
    %     \vlderivation{
    %     \vlin{\lr \mu}{}{\Gamma, \sigma(\tau) \seqar \sigma(\tau')}{
    %     \vlin{\rr \mu}{}{\Gamma, \sigma'(\tau, \sigma(\tau'))\seqar \sigma(\tau')}{
    %     \vliq{\sigma'(P,\sigma(\tau'))}{}{\Gamma, \sigma'(\tau, \sigma(\tau')) \seqar \sigma'(\tau', \sigma(\tau')) }{
    %     \vlhy{\Gamma, \tau \seqar \tau'}
    %     }
    %     }
    %     }
    %     }
    %     \]
    %     \todo{reference variations of iterator}
    %     where $\sigma'(P,\sigma(\tau'))$ is obtained from the IH for $\sigma'(P,Y)$ under substitution of $\sigma(\tau')$ for $Y$, cf.~Remark~\ref{rem:gen-identity+substitution}.
    %     Dually for $\rho(P)$.
    %     \item Similarly for greatest fixed points. 
    % \end{itemize}
\end{definition}

\begin{example}
[Post-fixed point]
\label{ex:post-fixed-in-muLJ}
It is implicit in the rules of $\muLJ$ that $\mu X \sigma(X)$ may be seen as the \emph{least} fixed point of $\sigma(\cdot)$, under a suitable semantics (e.g.\ later in Section~\ref{sec:totality}).
The $\rr \mu $ rule indicates that it is a pre-fixed point, while the $\lr \mu$ rule indicates that it is least among them. 
To see that it is also a post-fixed point $\small\vlinf{\mulunf}{}{\Gamma, \mu X \sigma(X)\seqar \tau}{\Gamma, \sigma(\mu X \sigma (X))\seqar \tau}$ we may use a derivation that mimics standard textbook-style proofs of Knaster-Tarski:
\[
\small
\vlderivation{
\vliin{\lr \mu}{}{\Gamma, \mu X \sigma(X) \seqar \tau}{
    \vliq{\sigma}{}{ \sigma(\sigma(\mu X \sigma(X)))\seqar \sigma(\mu X \sigma(X))}{
    \vlin{\rr \mu}{}{ \sigma(\mu X\sigma(X))\seqar \mu X \sigma(X)}{
    \vlin{\id}{}{ \sigma(\mu X \sigma(X))\seqar \sigma(\mu X \sigma(X))}{\vlhy{}}
    }
    }
}{
    \vlhy{\Gamma, \sigma(\mu X \sigma(X))\seqar \tau}
}
}
\]
Dually, we can derive  $\vlinf{\nurunf}{}{\Gamma \seqar \nu X \sigma(X)}{\Gamma \seqar \sigma(\nu X \sigma(X))}$ as follows: 
\[
\small
\vlderivation{
\vliin{\rr \nu}{}{\Gamma\seqar  \nu X \sigma(X) }
{
    \vlhy{\Gamma \seqar  \sigma(\nu X \sigma(X))}
}
{
    \vliq{\sigma}{}{\sigma(\nu X \sigma(X)) \seqar \sigma(\sigma(\nu X \sigma(X))) }{
    \vlin{\lr \nu}{}{ \nu X \sigma(X)\seqar \sigma(\nu X\sigma(X)) }{
    \vlin{\id}{}{ \sigma(\nu X \sigma(X))\seqar \sigma(\nu X \sigma(X))}{\vlhy{}}
    }
    }
}
}
\]
% for $\nu$-types as pre-fixed points.
% (see~\Cref{app:simple-types-with-fixed-points}, point~\ref{enum:app1}).
% $\vlinf{\nurunf}{}{\Gamma \seqar \nu X \tau(X)}{\Gamma \seqar \tau(\nu X \tau(X))}$. 
\end{example}

\subsection{Computing with derivations}
The underlying computational model for sequent calculi, with respect to the `proofs-as-programs' paradigm, is \emph{cut-reduction}.
In our case this follows a standard set of cut-reduction rules for the calculus $\LJ$.
For the fixed points, cut-reduction is inherited directly from the encoding of fixed points in system $\F$ that induces our rules, cf.~Remark~\ref{rem:mulj-frag-F}.
Following Baelde and Miller~\cite{BaeldeMiller07}, we give self-contained cut-reductions here:
\gianluca{I have put my/Farzad's definition of cut-elimination for lrmu vs rr mu. I commented the notion of application. Maybe some of the discussion above is now obsolete?}
\anupam{16-10-23: it might be worth giving the application-version of cut-reduction here too, at least as an example: it exemplifies how cut-reduction in $\muLJ$ specialises that of second-order logic.}


 \gianluca{Also, most of the following discussion above can be now found in the appendix. Anupam, please check if you are fine with it: NB: contractions above occur because we have $\Delta$ context in the step function. This seems to break cut-elimination in linear logic, hence why Baelde and Miller have context-free step function. 
        Note that it is on this point that the format of the rules must (crucially) differ from that given by F, where the presence of context in the step function does not compromise linearity.
        Under an appropriate semantics, e.g. in term calculi section, or even embedding into circular system, it makes no difference, but we might want to change the rule presentation globally. Update: in fact perhaps this is not best at all, since without a context in the step function it is not clear how to cut-free derive the iterator! it would be nice for the iterator to be cut-free derivable, not least for embedding target of realisability. Under suitable semantics, e.g. term calculi, all choices are equivalent (using cuts).}

% \begin{defn}
%     [Application and cut-reduction for fixed points]
%     \label{def:appl-and-cut-red-muLJ}
%     Given $P:\Gamma\seqar \mu X \sigma(X)$ and $Q:\Delta,\sigma(\rho) \seqar \rho$   we define the \emph{application} of $P$ to $Q$, written simply $PQ$, as a proof of $\Gamma ,\Delta \seqar \rho$ as follows:
%     append $\Delta$ to the LHS everwhere in $P$ and replace every descendant $\sigma'(\mu X \sigma(X))$ on the RHS by $\sigma'(\rho)$ \todo{need to justify format of descendants}  and locally repair the proof by:
%     \begin{itemize}
%         \item affected $\rr \mu$ steps are transformed as follows:
%         \[
%         \small
%         \vlinf{}{}{\Sigma  \seqar \mu X \sigma(X)}{\Sigma \seqar \sigma(\mu X\sigma(X))}
%         \  \mapsto \ 
%         \vlderivation{
%         \vliq{\contr}{}{\Sigma, \Delta \seqar \rho}{
%         \vliin{\cut}{}{\Sigma, \Delta, \Delta \seqar \rho}{
%             \vlhy{ \Sigma, \Delta \seqar \sigma(\rho)}
%             }{
%             \vltr{Q}{\Delta, \sigma(\rho)\seqar \rho}{\vlhy{\ }}{\vlhy{}}{\vlhy{\ }}
%             }
%         }
%         }
%         \]
        
%         \item an affected identity on $\sigma'(\mu X \sigma(X))$ is transformed to:
%         \[
%      \small
%         \hspace{-1.5em}
%         % \vlinf{\id}{}{\sigma'(\mu X \sigma(X))\seqar \sigma'(\mu X \sigma(X))}{}
%         % \quad \leadsto \quad
%         \vlderivation{
%         \vliq{\sigma'}{}{\Delta, \sigma'(\mu X \sigma(X))\seqar \sigma'(\rho)}{
%         \vlin{\lr \mu}{}{\Delta, \mu X\sigma(X) \seqar \rho}{
%         \vltr{Q}{\Delta, \sigma(\rho) \seqar \rho}{\vlhy{\ }}{\vlhy{}}{\vlhy{\ }}
%         }
%         }
%         }
%         \ \text{or}\ 
%         \vlderivation{
%         \vliq{\sigma'}{}{\Delta, \sigma'(\rho)\seqar \sigma'(\mu X \sigma(X))}{
%         \vlin{\lr \mu}{}{\Delta, \mu X\sigma(X) \seqar \rho}{
%         \vltr{Q}{\Delta, \sigma(\rho) \seqar \rho}{\vlhy{\ }}{\vlhy{}}{\vlhy{\ }}
%         }
%         }
%         }
%         \]
%         depending on whether $\sigma'(\cdot)$ is positve or negative, resp.
%     \end{itemize}
% Dually for greatest fixed points.
% \todo{this can go in an appendix probably, or maybe the realisability part.}

\begin{defn}
    [Cut-reduction for fixed points]
    \label{def:appl-and-cut-red-muLJ}
\emph{Cut-reduction} on $\muLJ$-derivations, written $\cutred$, is the smallest relation on derivations including all the usual cut-reductions of $\LJ$ and the reductions in~\Cref{fig:cut-reduction-mu-nu-muLJ-paper}. {As usual, we allow these reductions to be performed on sub-derivations. I.e.\ they are `context-closed'.} 
% 
% \[
% \small
% \begin{array}{rl}
% & \vlderivation{
% \vliin{\cut}{}{\Gamma, \Delta,\Sigma \seqar \tau}{
%     \vltr{\der}{\Gamma \seqar \mu X \sigma(X)}{\vlhy{\ }}{\vlhy{}}{\vlhy{\ }}
% }{
%     \vliin{\lr \mu}{}{\Delta, \Sigma, \mu X \sigma(X) \seqar \tau}{
%         \vltr{\deri}{\Delta,\sigma(\rho ) \seqar \rho}{\vlhy{\ }}{\vlhy{}}{\vlhy{\ }}
%     }{
%         \vltr{\derii}{\Sigma, \rho \seqar \tau}{\vlhy{\ }}{\vlhy{}}{\vlhy{\ }}
%     }
% }
% } \\
% \rewrite & 
% \vlderivation{
% \vliin{\cut}{}{\Gamma, \Delta,\Sigma \seqar \tau}{
%     \vltr{\der}{\Gamma \seqar \mu X \sigma(X)}{\vlhy{\ }}{\vlhy{}}{\vlhy{\ }}
% }{
%     \vliin{\lr \mu}{}{\Delta, \Sigma, \mu X \sigma(X) \seqar \tau}{
%         \vltr{\deri}{\Delta,\sigma(\rho ) \seqar \rho}{\vlhy{\ }}{\vlhy{}}{\vlhy{\ }}
%     }{
%         \vltr{\derii}{\Sigma, \rho \seqar \tau}{\vlhy{\ }}{\vlhy{}}{\vlhy{\ }}
%     }
% }
% }
% \end{array}
% \]
\end{defn}

 
\begin{figure}
    \centering
%     \[
%     \footnotesize
% %     \vlderivation{
% %     \vliin{\cut}{}{\Sigma, \Gamma,\Delta\seqar \tau}
% %     {
% %      \vlin{\rr \mu}{}{\Sigma \seqar \mu X \sigma(X)}{
% %      \vlhy{\Sigma \seqar \sigma(\mu X \sigma(X))}}
% %      }
% %      {
% %   \vliin{\lr \mu}{}{\Gamma,\Delta, \mu X \sigma(X)  \seqar \tau}
% %   {\vlhy{\Gamma, \sigma(\rho)\seqar \rho}}{\vlhy{\Delta, \rho \seqar \tau}}    
% %   }
% %     }
% % \!\quad \rightsquigarrow \! \quad
% \vlderivation{
% \vliq{\contr}{}{\Sigma, \Delta,  \Gamma\seqar \tau}
% {
% \vliin{\cut}{}{\Sigma,  \Delta, \Gamma, \Gamma \seqar \tau}
% {
% \vliin{\cut}{}{\Sigma, \Gamma, \Gamma \seqar \rho}
% {
% \vliin{\cut}{}{\Sigma, \Gamma\seqar \sigma(\rho)}
% {
% \vlhy{\Sigma \seqar \sigma(\mu X \sigma(X))}
% }
% {
% \vliq{\sigma}{}{\Gamma, \sigma(\mu X \sigma(X) ) \seqar \sigma(\rho)}
% {
%  \vlin{\lr \mu}{}{\Gamma, \mu X \sigma(X)  \seqar \rho}
%   {\vlhy{\Gamma, \sigma(\rho)\seqar \rho}}
%   % {\vlin{\id}{}{ \rho \seqar \rho}{\vlhy{}}}    
% }
% }
% }
% {\vlhy{\Gamma, \sigma(\rho)\seqar \rho}}
% }
% {
% \vlhy{\Delta, \rho \seqar \rho}
% }
% }
% }
% \]
% \[
%  \small
% \vlderivation{
% \vliin{\cut}{}{\Delta, \Gamma , \Sigma \seqar \tau }
% {
%  \vliin{\rr \nu }{}{\Delta,  \Gamma \seqar \nu X\sigma(X)}{\vlhy{\Delta \seqar  \rho }}{\vlhy{  \Gamma, \rho \seqar \sigma(\rho)}}
%  }
%  {
%  \vlin{\lr \nu }{}{\Sigma,  \nu X\sigma(X) \seqar \tau}{\vlhy{\Sigma , \sigma(\nu X \sigma(X))\seqar \tau}}
%  }
%  }
% \qquad \rightsquigarrow \qquad
% \vlderivation{
% \vliq{\contr}{}{\Delta , \Gamma, \Sigma \seqar \tau}
% {
% \vliin{\cut}{}{\Delta , \Gamma, \Gamma, \Sigma , \Delta \seqar \tau}
% {\vlhy{\Delta \seqar \rho}}
% {
% \vliin{\cut}{}{\Gamma, \Gamma, \Sigma , \rho \seqar \tau}
% {
% \vlhy{  \Gamma, \rho \seqar \sigma(\rho)}
% }
% {
% \vliin{\cut}{}{\Gamma ,  \Sigma, \sigma(\rho) \seqar \tau}
% {
% \vliq{\sigma }{}{\Gamma ,  \sigma(\rho) \seqar \sigma(\nu X\sigma(X))}
% {
% \vliin{\rr \nu }{}{\Gamma , \rho \seqar \nu X\sigma(X)}{\vlin{\id}{}{\rho \seqar  \rho}{\vlhy{} }}{\vlhy{  \Gamma, \rho \seqar \sigma(\rho)}}
% }
% }
% {\vlhy{\Sigma , \sigma(\nu X \sigma(X))\seqar \tau}}
% }
% }
% }
% }
% \]
% \begin{figure*}
    \centering
    \[
    \footnotesize
    \vlderivation{
    \vliin{\cut}{}{\Sigma, \Gamma,\Delta\seqar \tau}
    {
     \vlin{\rr \mu}{}{\Sigma \seqar \mu X \sigma(X)}{
     \vlhy{\Sigma \seqar \sigma(\mu X \sigma(X))}}
     }
     {
  \vliin{\lr \mu}{}{\Gamma,\Delta, \mu X \sigma(X)  \seqar \tau}
  {\vlhy{\Gamma, \sigma(\rho)\seqar \rho}}{\vlhy{\Delta, \rho \seqar \tau}}    
  }
    }
    \qquad
    \quad
    \rightsquigarrow
\]
\vspace{0.2cm}
\[
\footnotesize
\vlderivation{
\vliq{\contr}{}{\Sigma, \Delta,  \Gamma\seqar \tau}
{
\vliin{\cut}{}{\Sigma,  \Delta, \Gamma, \Gamma \seqar \tau}
{
\vliin{\cut}{}{\Sigma, \Gamma, \Gamma \seqar \rho}
{
\vliin{\cut}{}{\Sigma, \Gamma\seqar \sigma(\rho)}
{
\vlhy{\Sigma \seqar \sigma(\mu X \sigma(X))}
}
{
\vliq{\sigma}{}{\Gamma, \sigma(\mu X \sigma(X) ) \seqar \sigma(\rho)}
{
 \vliin{\lr \mu}{}{\Gamma, \mu X \sigma(X)  \seqar \rho}
  {\vlhy{\Gamma, \sigma(\rho)\seqar \rho}}{\vlin{\id}{}{ \rho \seqar \rho}{\vlhy{}}}    
}
}
}
{\vlhy{\Gamma, \sigma(\rho)\seqar \rho}}
}
{
\vlhy{\Delta, \rho \seqar \tau}
}
}
}
\]
\vspace{1cm}
\[
 \footnotesize
\vlderivation{
\vliin{\cut}{}{\Delta, \Gamma , \Sigma \seqar \tau }
{
 \vliin{\rr \nu }{}{\Delta,  \Gamma \seqar \nu X\sigma(X)}{\vlhy{\Delta \seqar  \rho }}{\vlhy{  \Gamma, \rho \seqar \sigma(\rho)}}
 }
 {
 \vlin{\lr \nu }{}{\Sigma,  \nu X\sigma(X) \seqar \tau}{\vlhy{\Sigma , \sigma(\nu X \sigma(X))\seqar \tau}}
 }
 }
 \qquad
 \quad
 \rightsquigarrow
\]
\vspace{0.2cm}
\[
\footnotesize
\vlderivation{
\vliq{\contr}{}{\Delta , \Gamma, \Sigma \seqar \tau}
{
\vliin{\cut}{}{\Delta , \Gamma, \Gamma, \Sigma , \Delta \seqar \tau}
{\vlhy{\Delta \seqar \rho}}
{
\vliin{\cut}{}{\Gamma, \Gamma, \Sigma , \rho \seqar \tau}
{
\vlhy{  \Gamma, \rho \seqar \sigma(\rho)}
}
{
\vliin{\cut}{}{\Gamma ,  \Sigma, \sigma(\rho) \seqar \tau}
{
\vliq{\sigma }{}{\Gamma ,  \sigma(\rho) \seqar \sigma(\nu X\sigma(X))}
{
\vliin{\rr \nu }{}{\Gamma , \rho \seqar \nu X\sigma(X)}{\vlin{\id}{}{\rho \seqar  \rho}{\vlhy{} }}{\vlhy{  \Gamma, \rho \seqar \sigma(\rho)}}
}
}
{\vlhy{\Sigma , \sigma(\nu X \sigma(X))\seqar \tau}}
}
}
}
}
\]
    \caption{Cut-reduction rules for $\mu$ and $\nu$ in $\muLJ$.}
    \label{fig:cut-reduction-mu-nu-muLJ}
    % \caption{Reduction of a key cut between $\rr \mu$ and $\lr \mu$ in $\muLJ$. \todonew{add LHS of the reduction.}}
    \label{fig:cut-reduction-mu-nu-muLJ-paper}
\end{figure}



% Let us henceforth write $P\cutred Q$ if $P$ reduces to $Q$ under a cut-reduction step, and similarly $\cutred^*$ and $\cutconv$ for the reflexive transitive and reflexive symmetric transitive closures of $\cutred$, respectively.

When speaking of (subsets of) $\muLJ$ as a computational model, we always mean with respect to the relation $\cutred^*$ unless otherwise stated.
More precisely:

\begin{defn}[Representability in $\muLJ$]\label{defn:numerals-representability}
We define the \emph{type of natural numbers} as $\N\dfn \mu X(1+ X)$. 
We also define the \emph{numeral} $\cod{n}: \N$  by induction on $n \in \Nat$: 
\[
 \small
\begin{array}{rclrcl}
 \cod 0    & \dfn&
     \vlderivation{
     \vlin{\rr \mu}{}{\seqar \N}{\vlin{ {\rr +}}{}{\seqar \unit + \N}{\uaxiom }}
     }
     &\qquad 
    \cod {n+1} & \dfn & 
    \vlderivation{
    \vlin{\rr \mu }{}{\seqar\N}{\vlin{{\rr +}}{}{\seqar \unit + \N}{
    \vliq{\numeral n}{}{\seqar \N}{\vlhy{}}
    }
    }
    }
\end{array}
\]    
We say that a {(possibly partial)} function $f: \Nat \times \overset{k}{\ldots} \times \Nat \to \Nat$ is \emph{representable} in $\muLJ$ \todonew{or any other system? maybe mention that muLJ is known to be total} if there is  a $\muLJ$-derivation $\der_f: \N , \overset{k}{\ldots} ,    \N \seqar \N  $ s.t., for any $n_1, \ldots, n_k \in \Nat$, the derivation,
\[\small
\vlderivation{
\vliin{\cut}{}{\seqar \N}{\vliq{\numeral n_k}{}{\seqar \N}{\vlhy{}}}{
\vliin{\cut}{}{\vdots}{ \vliq{\numeral n_2}{}{\seqar \N}{\vlhy{}} }{
\vliin{\cut}{}{\N,\overset{k-1}{\ldots},  \N\seqar \N}{ \vliq{\numeral n_k}{}{\seqar \N}{\vlhy{}} }{\vltr{\der_f}{\N,\overset{k}{\ldots},  \N\seqar \N}{\vlhy{\ }}{\vlhy{\ }}{\vlhy{\ }} }
}
}
}
% \qquad
% \cutelim^*
% \qquad 
% \vlderivation{
% \vldr{\cod{f(n_1, \ldots, n_k)}}{\seqar \N}
% }
\]
reduces under $\cutred^*$ to the numeral $\numeral{f(n_1, \dots, n_k)}$, whenever it is defined (otherwise it reduces to no numeral).
In this case we say that $\der_f$ \emph{represents} $f$ in $\muLJ$.
\end{defn}

\begin{figure}
    \[
    \vlinf{\Nzero}{}{\seqar \N}{}
    \qquad 
    \vlinf{\Nsucc}{}{\Gamma \seqar \N}{\Gamma \seqar \N}
    \qquad
    \vliiinf{\lr \N}{}{\Gamma,\Delta, \N \seqar \tau}{\Gamma \seqar \sigma}{\Gamma, \sigma \seqar \sigma}{\Delta, \sigma \seqar \tau}
    % \quad
    % \vliinf{\cnd}{}{\Gamma, \N \seqar \tau}{\Gamma \seqar \tau}{\Gamma, \N\seqar \tau}
    \]
    \caption{Native rules for $\N$ in $\muLJ$.}
    \label{fig:native-n-rules-mulj}
\end{figure}
\begin{example}
    [Native rules for natural number computation]
    `Native' rules for type $\N$ in $\muLJ$ are given in Figure~\ref{fig:native-n-rules-mulj}, all routinely derivable in $\muLJ$, as in Figure~\ref{fig:deriving-nat-rules-in-muLJ}.
    The corresponding `native' cut-reductions, derivable using $\cutred$ cf.~\Cref{fig:cut-reduction-mu-nu-muLJ}, are also routine.
    We shall examine this further in \Cref{sec:term-calculi}.
    \todonew{give cut reductions?}
     \gianluca{Say that recursor is iteration: we can get proper recursion using pairing.}
    Note that, from here we can recover the usual recursor of system $\T$, as shown formally by Clairambault~\cite{Clairambault13interleaving}.
\end{example}



% \subsection{Further examples}

\begin{example}\label{exmp:examples-of-derivations}
The least and greatest fixed point operators $\mu$ and $\nu$ allow us to encode inductive data (natural numbers, lists, etc) and coinductive data (streams, infinite trees, etc). We have already seen the encoding of natural numbers. The type of lists and streams (both over natural numbers) can be represented by, respectively,  $L\dfn \mu X (\unit + (\N \times X))$ and $S= \nu X (\N \times X)$.
Figure~\ref{fig:examples-derivations}, left-to-right, shows the encoding of  $\varepsilon$ and $n::l$ (i.e., the empty list and the operation appending  a natural number to a list). 
Figure~\ref{fig:list+stream-to-stream-muLJ} shows the encoding of a concatenation of a list and a stream into a stream (by recursion over the list with the invariant $S \to S$).\footnote{Both examples were originally given for $\muMALL$ in \cite{Doumane17thesis}.}
\end{example}
% \anupam{could give some example of stream computation here, since we barely deal with greatest fixed points in the formalities}

\begin{figure}
    \centering
     \[
\Nzero\ \dfn \ 
\vlderivation{
\vlin{\rr \mu}{}{\seqar \nat}{
\vlin{\rr +^0}{}{\seqar 1 + \nat}{
\vlin{\rr 1 }{}{\seqar 1}{\vlhy{}}
}
}
}
\qquad
 \Nsucc\ \dfn \ 
\vlderivation{
\vlin{\rr \mu}{}{\Gamma \seqar N}{
\vlin{\rr {+^1}}{}{\Gamma\seqar \unit + \N}
{
\vlhy{\Gamma \seqar \N}
}
}
}
\qquad 
\lr N \ \dfn \ \vlderivation{
\vliin{\lr\mu}{}{\Gamma, \Delta, \N \seqar \tau}
{
\vliin{\lr +}{}{\Gamma, 1 + \sigma \seqar \sigma }
{
\vlin{\lr \unit}{}{\Gamma, \unit \seqar \sigma }{\vlhy{\Gamma \seqar \sigma}}
}
{
\vlhy{\Gamma, \sigma \seqar \sigma}
}
}
{
\vlhy{\Delta, \sigma \seqar  \tau}
}
}
    \]
    \caption{Constructing and destructing natural numbers in $\muLJ$.}
    \label{fig:deriving-nat-rules-in-muLJ}
\end{figure}

\begin{figure}
\[
\small
\vlderivation{
\vlin{\rr \mu}{}{\seqar L}
{
\vlin{\rr {+^0}}{}{\seqar\unit + (\N \times L)}
{
\vlin{\rr \unit}{}{\seqar\unit}{\vlhy{}}
}
}
}
\qquad 
\vlderivation{
\vlin{\rr \mu}{}{\seqar L}
{
\vlin{\rr{+^1}}{}{\seqar \unit +(\N \times L) }
{
\vliin{\times}{}{\seqar \N \times L}
{
\vltr{\cod{n}}{\seqar \N}{\vlhy{}}{\vlhy{}}{\vlhy{}}
}
{
\vltr{\underline{l}}{\seqar L}{\vlhy{}}{\vlhy{}}{\vlhy{}}
}
}
}
}
\]
        \caption{Constructing lists in $\muLJ$.}
    \label{fig:examples-derivations}
\end{figure}

\begin{figure}
    \centering
\[
\small
\vlderivation{
\vlin{\lr \nu}{}{L \seqar S \to S}
{
\vliin{\lr + }{}{\unit +(N \times (S \to S))  \seqar S \to S}
{
\vlin{\lr \unit}{}{\unit \seqar S \to S}
{
\vlin{\rr \to }{}{ \seqar S \to  S}
{
\vlin{\id }{}{S \seqar S}{\vlhy{}}
}
}
}
{
\vlin{\lr \times}{}{N \times (S \to S) \seqar  S \to S}
{
\vlin{\rr \to}{}{N, S \to S \seqar  S \to S}
{
\vlin{\rr \nu}{}{N, S \to S, S \seqar S}
{
\vliin{\rr \times}{}{N, S \to S, S \seqar N \times S}
{
\vliin{\lr \to}{}{S \to S, S \seqar S}
{
\vlin{\id}{}{S \seqar S}{\vlhy{}}
}
{
\vlin{\id}{}{S \seqar S}{\vlhy{}}
}
}
{
\vlin{\id}{}{N \seqar N}{\vlhy{}}
}
}
}
}
}
}
}
\]    \caption{Concatenation of a list and a stream in $\muLJ$.}    \label{fig:list+stream-to-stream-muLJ}
\end{figure}
